\section{Алгоритм поиска наиболее важного объекта (НВО)}

Работа систем АКК и АЭТ напрямую связана с конкретными транспортными средствами. Для АКК это целевое ТС, то есть ТС, расстояние до которого контролируется. Для АЭТ это НВО могут быть как ТС, так и люди. Таким образом, необходимо разработать алгоритм по определению НВО среди всех ослеживаемых объектов, которые фиксируются алгоритмом DeepSORT. 

Алгоритм по определению НВО представлен на рисунке~\ref{fig:mio}

Функция фильтрует объекты на основе их скорости и положения, исключая статические объекты и объекты, движущиеся навстречу, а также объекты, находящиеся вне заданных зон~\cite{Ziegler2014}.

\img[!h]{fig:mio}{mio_algorithm.jpg}{Алгоритм по поиску НВО}{1}

\subsection{Программная реализация алгоритма по поиску НВО}

Программная реализация данного алгоритма представлена на языке MATLAB. Это обусловлено основной работой по моделированию и проверке алгоритма в ПО MATLAB, а также облегчило его реализацию. 

\textbf{Входные данные}

Функция принимает три входных аргумента:

\lstinline:MIO_data_empty: — структура, представляющая пустую (инициализированную) выходную шину данных НВО. Она содержит поле \lstinline:NumDetections: (количество обнаруженных НВО, тип \lstinline:uint8:) и массив структур \lstinline:Detection:, каждая из которых описывает характеристики объекта: \lstinline:TargetIndex: (индекс объекта, тип \lstinline:uint8:), \lstinline:x: (координата по оси X, тип \lstinline:single:), \lstinline::{y} (координата по оси Y, тип \lstinline:single:), \lstinline:vx: (скорость по оси X, тип \lstinline:single:) и \lstinline:vy: (скорость по оси Y, тип \lstinline:single:).

\lstinline:allRadars_objects_data: — структура, содержащая данные всех объектов, обнаруженных радаром и алгоритмом DeepSORT~\cite{Feng_2021}. Она включает массив структур \lstinline:Detection:, каждая из которых содержит поля \lstinline:TargetIndex:, \lstinline:x:, \lstinline:y:, \lstinline:vx: и \lstinline:vy:, аналогичные полям в \lstinline:MIO_data_empty:.

\lstinline:CAN_body: — структура, содержащая данные из CAN-шины транспортного средства. Она включает поле \lstinline:Vehicle.VEH_SPEED: (скорость транспортного средства, тип \lstinline:single:).

\textbf{Выходные данные}

Функция возвращает структуру \lstinline:MIO_data:, которая содержит данные о наиболее опасных объектах (MIO). Структура включает поле \lstinline:NumDetections: (количество обнаруженных MIO, тип \lstinline:uint8:) и массив структур \lstinline:Detection:, каждая из которых описывает характеристики объекта: \lstinline:TargetIndex:, \lstinline:x:, \lstinline:y:, \lstinline:vx: и \lstinline:vy:.

\textbf{Алгоритм работы:}

\begin{enumerate}

    \item Инициализация:

    \begin{itemize}

        \item Создается массив \lstinline:MIO_index: для хранения индексов наиболее важных объектов.

        \item Инициализируется счетчик \lstinline:NumMIO: для подсчета количества НВО.

    \end{itemize}

    \item Фильтрация объектов

    Для каждого объекта, обнаруженного радаром и алгоритмом DeepSORT, проверяется:

    \begin{itemize}

        \item Наличие объекта (индекс \lstinline:TargetIndex: не равен 0)
        
        \item Относительная скорость объекта (\lstinline:vx: и \lstinline:vy:)

        \item Положение объекта в заданных зонах (левый, правый и продольный коридоры)

    \end{itemize}

    Объекты, которые двигаются навстречу (\lstinline:vx: <= 0), имеют скорость ниже порога (\lstinline:v_target_rel <= BSM_static_objects_velocity_threshold:) и находятся вне заданных зон исключаются из рассмотрения.


    \item Определение НВО:

    \begin{itemize}

        \item Если объект удовлетворяет всем условиям (находится в заданных зонах и не исключен), он добавляется в список НВО

        \item Индекс объекта сохраняется в массиве \lstinline:MIO_index:

        \item Увеличивается счетчик \lstinline:NumMIO:.

    \end{itemize}

    \item Формирование выходной структуры:

    \begin{itemize}

        \item В структуру \lstinline:MIO_data: записываются данные о каждом объекте, включая его индекс, координаты и скорость

        \item В поле \lstinline:NumDetections: записывается общее количество НВО.

    \end{itemize}

\end{enumerate}

\textbf{Калибровочные параметры}

Программа использует глобальные калибровочные параметры, которые определяют границы зон и пороговые значения:

\begin{itemize}

    \item \lstinline:BSM_y_left_left_bound:, \lstinline:BSM_y_left_right_bound: — границы левой зоны по оси Y.

    \item \lstinline:BSM_y_right_left_bound:, \lstinline:BSM_y_right_right_bound: — границы правой зоны по оси Y.

    \item \lstinline:BSM_x_upper_bound:, \lstinline:BSM_x_lower_bound: — границы продольной зоны по оси X.

    \item \lstinline:BSM_static_objects_velocity_threshold: — порог скорости для исключения статических объектов.

    \item \lstinline:BSM_relative_static_objects_velocity_threshold: — порог относительной скорости для исключения статических объектов.

\end{itemize}

Так как алгоритм представляет из себя универсальную функцию, калибровочные параметры задаются под конкретное ТС и сферу его применения. 

\subsection{Связь алгоритма поиска НВО с DeepSORT}

Алгоритм поиска наиболее важного объекта (НВО) напрямую использует данные, предоставляемые алгоритмом DeepSORT, который отслеживает объекты на основе их положения, скорости и визуальных признаков, обеспечивая непрерывное обновление треков в реальном времени. Алгоритм НВО анализирует информацию, полученную после алгоритма отслеживания и данных с радаров, затем выделяет объекты, представляющие наибольшую угрозу или важность для безопасности движения.